\section{Introduction}
\label{sec:introduction}

Modern tensor computation frameworks face a fundamental tension between
performance optimization and execution correctness. Systems such as
TVM~\cite{chen2018tvm}, XLA~\cite{xla2017}, TensorRT~\cite{tensorrt2020},
and ONNX Runtime~\cite{onnxruntime2019} prioritize throughput by
aggressively applying optimizations---operator fusion, layout
transformation, quantization, kernel autotuning---without formally
verifying that the optimized execution plan preserves the numerical
semantics of the original computation. The gap between optimization and
correctness widens further under backend heterogeneity: a plan that is
valid on one accelerator may silently violate shape, layout, memory, or
numerical constraints on another.

This is not a hypothetical concern. During the empirical study reported
in Section~\ref{subsec:empirical}, an off-the-shelf release of a widely
deployed tensor compiler produced a ResNet-50 module whose output
diverged from four independent reference implementations by a relative
$L_2$ distance greater than $1.2$---changing the top-1
prediction---while compiling without a single warning. Fast, silent, and
wrong is precisely the failure mode this paper addresses.

\gate{} (Gwen Algorithm for Tensor Execution) is a
\emph{correctness-driven execution algorithm} built on one foundational
insight:

\begin{quote}
\itshape ``Numbers are not considered trustworthy until logical
constraints validate them.''
\end{quote}

A numerical result produced by a tensor computation is meaningless if
the execution plan that generated it violated one or more semantic
constraints. \gate{} inverts the traditional execution paradigm by
placing constraint validation \emph{before} cost optimization: every
plan considered by the optimizer has already been proven valid, which
eliminates the possibility of selecting a fast but incorrect execution
strategy.

\subsection{What \gate{} Is---and Is Not}

\gate{}'s position in the tensor execution ecosystem is easily
misunderstood, so we state it precisely. \gate{} is \textbf{not} a
neural-network framework (it defines no layers, losses, or training
loops), \textbf{not} a compiler (it parses no source and emits no
machine code), \textbf{not} an optimizer (it applies no loop
transformations or fusion itself), \textbf{not} a scheduler, and
\textbf{not} a tensor library. \gate{} \emph{is} the protocol layer that
sits above all of these components: it accepts candidate execution plans
as input, validates them against a composable set of logical
constraints, evaluates their cost under a policy-driven objective, and
selects the optimal valid plan for dispatch. \gate{} defines the
\emph{protocol} by which optimization and correctness interact, not the
implementation of either.

\subsection{Contributions}

This paper makes the following contributions:

\begin{enumerate}[leftmargin=2em]
  \item \textbf{Algorithm.} A four-phase execution algorithm
    (generate--validate--evaluate--dispatch) in which a constraint
    engine enforcing seven composable constraint types---shape, tensor
    layout, memory, backend capability, numerical error, determinism,
    and safety---runs strictly before any cost evaluation
    (Sections~\ref{sec:algorithm}--\ref{sec:constraint-engine}).
  \item \textbf{Formalization.} Mathematical definitions of execution
    plans, constraints, validators, the valid execution set, metric
    vectors, weight vectors, and the cost function, together with proofs
    of constraint soundness, dispatch correctness, and cost optimality
    (Sections~\ref{sec:math} and~\ref{sec:correctness}).
  \item \textbf{Complexity analysis.} Constraint validation adds
    $O(knp)$ overhead for $k$ constraints and $p$ candidate plans of $n$
    operations, with early-exit reducing the effective cost by
    $2.5$--$3.5\times$ in practice (Section~\ref{sec:complexity}).
  \item \textbf{Implementation.} A four-crate Rust implementation
    (glcore, glproc, glcuda, glvulkan), plus two independent,
    dependency-free reference implementations used for empirical
    validation (Section~\ref{sec:implementation}).
  \item \textbf{Evaluation.} An analytical study across ten models and
    five competing frameworks (Section~\ref{sec:evaluation}), now
    complemented by an end-to-end \emph{empirical} proof-of-concept on
    commodity CPU and real TPU~v5e hardware
    (Section~\ref{subsec:empirical}), in which the full 7-constraint
    validation and plan selection measure at $1.4$--$8.2\,\mu$s per
    policy---seven orders of magnitude below one inference---and in
    which the numerical-error constraint catches a real miscompilation
    in a production compiler.
\end{enumerate}
